\section{补充结果、软件与复跑说明}\label{app:supplement}
\subsection{材料组织与结果索引}
本论文工程的\texttt{evidence/claim\_index.json}登记正文论断、源文件哈希和图表锚点。原始研究表格保留在各问\texttt{round3/results/}，论文图及其绘图脚本位于论文工程\texttt{figures/}和\texttt{scripts/}。每个模型、机制、情景和几何版本沿用原标识，程序附录对应论文实际采用的自编计算与后处理代码。

\subsection{软件、版本与命令}
资料处理采用Python 3.12、NumPy、pandas、SciPy、matplotlib及按用途锁定的统计/地理环境；原共享环境含44个锁定包。传播计算采用本地Linux容器中的r.avaflow 4.0 Revision 7、D-Claw 1.0.0及Clawpack 5.11.0。r.avaflow资源收尾派生修复只涉及45处资源清理调用，有效共同配置的70项ASCII输出逐字一致，物理方程沿用固定源码。原版与派生镜像、每run实际源码/构建及配置快照分别登记。

论文的复跑从已有结果开始，依次执行证据/来源核对、论文绘图、参考文献排序、程序附录构建和XeLaTeX编译。完整命令在工程\texttt{README.md}。历史模型的重放命令见各问\texttt{REPRODUCE.md}，其输入版本和缓存索引与论文证据表关联。

\begin{table}[htbp]\centering\small
\caption{研究与论文重放入口。命令在仓库根执行，各路径指向相应版本的说明。}
\begin{tabularx}{\linewidth}{p{2.3cm}Y}
\toprule
模块 & 入口\\\midrule
Q1因素识别 & \texttt{track\_b/research/q1/round3/REPRODUCE.md}\\
Q2机制及实验 & \texttt{track\_b/research/q2/round3/REPRODUCE.md}\\
Q3时间预测 & \texttt{track\_b/research/q3/round3/REPRODUCE.md}\\
Q4传播模型 & \texttt{track\_b/research/q4/round3/REPRODUCE.md}\\
论文集成 & \texttt{track\_b/paper/README.md}\\\bottomrule
\end{tabularx}\end{table}

\subsection{数据与模型配置的层级}
气象排序同时保留主质量轨与探索质量轨；结构水热计算以共同256参数设计、四机制及两强度政策组织；时间预测将历史信息、假设延迟和理想天气输入分别登记；传播给出几何、释放体积、组分、密度、摩阻、孔压、格网、时间步及目标定义。详细参数表、分组评价和每次运行字段由证据索引链接原表，供逐项核对。

\subsection{结构水热共同参数设计}\label{app:parameter-table}
下表给出空间填充设计的19项参数上下限及取样尺度。四机制共享256个参数点，两强度政策形成2048个设计情景；参数设计与概率先验分别登记。实际各参数点见相应参数表。
\begingroup\small\begin{longtable}{p{3.8cm}p{2.2cm}rrp{1.6cm}}
\caption{共同参数设计范围。}\label{tab:app-parameters}\\
\toprule 参数 & 单位 & 下限 & 上限 & 取样尺度\\\midrule\endfirsthead
\toprule 参数 & 单位 & 下限 & 上限 & 取样尺度\\\midrule\endhead
结构面倾角 & $^\circ$ & 30 & 60 & 线性\\
法向厚度 & m & 10 & 60 & 线性\\
岩体体积分数 & 1 & 0.2 & 0.8 & 线性\\
岩桥承载面积比例 & 1 & 0 & 0.25 & 线性\\
岩桥有效黏聚力 & kPa & 300 & 3000 & 对数\\
岩面摩擦系数 & 1 & 0.45 & 0.85 & 线性\\
初始裂隙温度 & $^\circ$C & -6 & -0.8 & 线性\\
热扩散率 & m$^2$/s & 9e-07 & 2.1e-06 & 线性\\
热路径深度 & m & 0.5 & 20 & 对数\\
站点至源区温度偏移 & $^\circ$C & -18 & -6 & 线性\\
储水量/水头比 & 1 & 0.002 & 0.08 & 对数\\
排水时间常数 & h & 3 & 240 & 对数\\
最大连通水头 & m & 2 & 40 & 对数\\
有效入渗倍率 & 1 & 0.25 & 3 & 对数\\
度日融化系数 & mm/($^\circ$C$\cdot$d) & 0.5 & 6 & 线性\\
损伤速率系数 & h$^{-1}$ & 1e-07 & 0.0002 & 对数\\
应力比幂次 & 1 & 2 & 6 & 线性\\
初始损伤 & 1 & 0.05 & 0.7 & 线性\\
损伤排水耦合系数 & 1 & 0 & 3 & 线性\\
\bottomrule\end{longtable}\endgroup
\subsection{解析热机制分类与裕度补充}
分类按相应初态与首次过阈定义，在同一设计表内汇总；初始越阈与未来越阈分列。
\begin{table}[htbp]\centering\small\caption{最终解析热版本的机制分类与最小稳定裕度中位数。}\label{tab:app-mechanisms}
\begin{tabular}{llrrrr}\toprule 机制 & 强度政策 & 未来越阈 & 初始越阈 & 时域内保持 & 裕度中位数\\\midrule
S & 线性延伸 & 4 & 31 & 221 & 0.6901\\
S+T & 线性延伸 & 28 & 31 & 197 & 0.4853\\
S+H & 线性延伸 & 6 & 31 & 219 & 0.6623\\
S+T+H & 线性延伸 & 31 & 31 & 194 & 0.4718\\
S & 暖端岩接触 & 4 & 31 & 221 & 0.6901\\
S+T & 暖端岩接触 & 27 & 31 & 198 & 0.4853\\
S+H & 暖端岩接触 & 6 & 31 & 219 & 0.6623\\
S+T+H & 暖端岩接触 & 31 & 31 & 194 & 0.4718\\
\bottomrule\end{tabular}\end{table}


地形派生成果使用Copernicus WorldDEM-30（DLR e.V. 2010—2014、Airbus Defence and Space GmbH 2014—2018，由欧盟与ESA通过COPERNICUS提供）；光学图件包含经修改的Copernicus Sentinel 2026数据。对象与河网资料署名为OpenStreetMap contributors，依ODbL登记。
